\section{Introduction}
Ragging and cyberbullying remain persistent problems in Indian higher education despite regulation~\cite{ugc2009}.
SafeVoice is a victim-initiated reporting tool: a student pastes a message, a DistilBERT toxicity classifier and a zero-shot
DistilBART context classifier produce a composite severity score in $[0,1]$, LIME explains the decision~\cite{ribeiro2016,sanh2019},
and \emph{critical} cases (score $\ge 0.75$) trigger an immediate counsellor alert.
The NLP pipeline decides \emph{how urgent} a report is; it says nothing about whether the alert \emph{arrives}.
In a hostel the delivery path is the building's Wi-Fi and the Internet uplink behind it, and these are least reliable exactly when many
people are online (nights, exams) or when an incident causes a surge of traffic.

This paper studies the transport layer underneath SafeVoice. We ask three questions:
(Q1) if alerts travel over an infrastructure-independent multi-hop ad-hoc mesh formed by the students' phones, does tagging them with
the NLP severity and prioritising them by severity improve delivery of the most severe alerts compared with the same mesh without priority?
(Q2) how does the mesh compare with a conventional infrastructure Wi-Fi deployment, both when the infrastructure is well provisioned
and when its shared uplink is congested? (Q3) which priority mechanism (802.11e EDCA, a severity-aware queue discipline, or both) is responsible
for any gain, and how do the systems behave under node or access-point failure and during a sudden crisis burst?

\paragraph{Contributions.}
(i) A single, tested mapping from the application's severity score to network classes (DSCP/802.11e access category/queue band).
(ii) An ns-3 study framework with explicit fairness controls: one scenario-independent workload file whose hash is logged per run
(identical positions, mobility, traffic and failures in all scenarios), identical PHY, queue sizes and rate settings, and a push-out
priority queue discipline with a unit test.
(iii) A validation record. Four modelling artefacts that would each have produced misleading results were detected by sanity gates and fixed
(Section~\ref{sec:validation}); we report them because they are easy to reproduce in other ns-3 studies.
(iv) An honest evaluation, including conditions where the proposed design does not help and where infrastructure Wi-Fi is simply better.
All code, workloads and raw per-packet logs are released with the paper.

\section{Background and Related Work}
\textbf{Ad-hoc routing.} AODV~\cite{rfc3561} discovers routes on demand by flooding route requests; its behaviour in dense networks is
dominated by control traffic and by the broadcast-storm effect~\cite{ni1999}. Comparative studies of MANET protocols~\cite{broch1998,das2000}
motivate our choice of a reactive protocol for a mostly static, partly mobile hostel. Link quality matters: shortest-hop routes found over
weak links perform badly~\cite{decouto2003}; we observed exactly this effect (Section~\ref{sec:validation}).
\textbf{Service differentiation.} IEEE 802.11e EDCA gives four access categories (voice, video, best effort, background) different
contention parameters and transmit-opportunity limits~\cite{ieee80211,mangold2003}; the saturation behaviour of the underlying
DCF is analysed in~\cite{bianchi2000}. EDCA differentiates \emph{channel access}; it does not by itself decide which packet is dropped when
a buffer fills, which is why we also evaluate a queue-level mechanism.
\textbf{Simulation validity.} Wireless simulators embed assumptions (propagation, interference, capture) that can change conclusions~\cite{kotz2004};
we therefore state every PHY/propagation assumption and list threats to validity explicitly. We use ns-3~\cite{riley2010}.

\section{System Model}
\subsection{From text severity to network class}
The application computes a severity score $s\in[0,1]$ (50\% model probability, 50\% weighted context flags) and maps it to
CRITICAL ($s\ge0.75$), HIGH ($s\ge0.50$), MODERATE ($s\ge0.25$) or LOW. The network study uses three SOS classes
(Table~\ref{tab:map}); \texttt{tools/severity\_bridge.py} is the single definition of the mapping and a unit test parses the application's
configuration to ensure that the thresholds cannot drift apart.
\begin{table}[t]\centering\caption{Severity-to-network mapping used in scenario A.}\label{tab:map}
\small\begin{tabular}{@{}lcccc@{}}\toprule
App risk level & SOS class & DSCP & 802.11e AC & queue band\\\midrule
CRITICAL ($\ge0.75$) & 3 & CS6 & AC\_VO & 3\\
HIGH ($\ge0.50$) & 2 & CS5 & AC\_VI & 2\\
MODERATE/LOW & 1 & 0 & AC\_BE & 1\\
other traffic & -- & 0 & AC\_BE & 0\\\bottomrule\end{tabular}\end{table}

\subsection{Scenarios}
\textbf{A (proposed)}: 802.11n ad-hoc mesh, AODV, severity priority (EDCA mapping of Table~\ref{tab:map} \emph{and} a severity-aware queue discipline).
\textbf{B (fair baseline)}: identical mesh and routing, no priority: all traffic is best effort in a FIFO queue.
\textbf{C (current practice, provisioned)}: infrastructure Wi-Fi: 8 access points (two per floor), each on its own orthogonal 20~MHz channel,
each station statically bound to the AP with the strongest computed signal, APs bridged through a gigabit wired core to a campus server that
hosts the SOS sink; SOS is ordinary best-effort traffic.
\textbf{C-up (current practice, congested uplink)}: as C, but the shared link between the campus core and the server is capacity limited
to $0.3\times$ the mesh saturation capacity, so that it saturates at the same offered load as the mesh (30\% of background flows are gateway-bound).
C-up models the mechanism by which hostel Wi-Fi congests an Internet-dependent alert: a shared bottleneck, not radio capacity.

\section{Simulation Methodology}\label{sec:method}
\subsection{Environment and workload}
All experiments use ns-3.44 (PyPI \texttt{ns3==3.44.post0} libraries; the planned 3.40 source tarball was not reachable from our build
environment) compiled with g++~13.3, C++20, on a 2-vCPU Linux container. The hostel has 4 floors of $40\times15$~m with a central corridor
and 12 rooms per floor (48 rooms). $n$ student phones (default $n=50$; $n\in\{25,75,100\}$ in the sweeps) plus one security desk
(ground-floor corridor) are placed as follows: 80\% of students are static at uniformly random positions in random rooms; 20\% are mobile,
following a random-waypoint model inside their floor's corridor (speed $U(1,1.5)$~m/s, pause $U(0,10)$~s).

\paragraph{Propagation and PHY (identical in A, B, C).} Log-distance path loss (exponent 3.0, 46.68~dB at 1~m) plus a deterministic symmetric
loss of 5~dB per wall and 13~dB per floor slab (a \texttt{PropagationLossModel} subclass: walls are counted analytically along the
straight line; floors from the height difference). 802.11n, 5~GHz, 20~MHz, one spatial stream, fixed HT-MCS3 (26~Mb/s) with
ConstantRate, 16.02~dBm transmit power, receiver sensitivity $-92$~dBm, ns-3's NIST error model, default A-MPDU aggregation. Broadcast frames are sent at the
data MCS. Ad-hoc scenarios share one channel (36); scenario C uses channels 36--64, one per AP.

\paragraph{Queues (identical everywhere).} Every station uses an explicit multi-queue root (\texttt{mq}) with four children of 100 packets,
one per access category (we do not rely on ns-3's default FQ-CoDel), and every Wi-Fi MAC queue is limited to 100 packets and 5~s
lifetime (ns-3's default 500~ms lifetime would silently expire packets and would be mislabelled as congestion). In A the children are
severity push-out queue discs; in B and C they are FIFO.

\paragraph{SOS traffic.} 160-byte UDP datagrams, Poisson arrivals at 3/s network-wide from uniformly random students to the desk; severity 3/2/1 with
probabilities 10/30/60\%; no application retries. At $t=30$~s a crisis burst adds 20 SOS from 20 distinct students within 5~s
(forced mix 5 critical/7 high/8 low). Each packet carries a header with sequence, transmit time, severity and source; all per-packet
statistics are computed from it. The first 10~s are a warm-up (every station pings the desk once so routes/associations exist; excluded
from statistics), followed by a 60~s measurement window (90~s with a failure at $t=50$~s in the robustness experiment); no SOS is generated in
the last 5~s.

\paragraph{Background traffic.} Video-like CBR flows (0.25--0.5~Mb/s, 1200~B) carry 70\% of the offered load and web-like on/off flows
(0.5~Mb/s while ON; ON/OFF exponential with means 1/3~s, 1000~B) 30\%; 30\% of flows go to the gateway (desk or server), the rest to random peers.
Flow rates are scaled down from the video rates in our original plan because a 1--2~Mb/s stream saturates a multi-hop 20~MHz mesh by itself.
\emph{Load levels} are defined as 10/40/80/120\% of the saturation capacity of the ad-hoc mesh, obtained from a calibration sweep
(Section~\ref{sec:calib}).

\subsection{Priority mechanisms}
\emph{EDCA level}: the sender sets the IP ToS from the SOS class (Table~\ref{tab:map}); ns-3 maps DSCP$\gg3$ to the user priority and
hence the access category at every hop, because forwarding preserves the IP header.
\emph{Queue level}: \texttt{SvPriorityQueueDisc} has four severity bands, strict-priority dequeue, FIFO within a band, and \emph{push-out}:
when the disc is full and a strictly lower non-empty band exists, the oldest packet of the lowest non-empty band is dropped and the arrival is accepted;
otherwise the arrival is dropped. (ns-3's \texttt{PrioQueueDisc} does not push out.) The severity is read from the UDP destination port
(7001--7003), independently of the ToS, so the queue mechanism can be evaluated without EDCA. A unit test checks dequeue order and
push-out accounting. Because \texttt{mq} gives each access category its own child, the queue discipline only reorders packets \emph{within} an
access category; in the full design that is class~1 versus background inside AC\_BE. The ablation therefore uses EDCA only (FIFO queues),
queue only (all traffic in AC\_BE, severity from the port) and both.

\subsection{Load calibration}\label{sec:calib}
A pilot (scenario B, $n=50$, run~1, 60~s window) offered 1--12~Mb/s of background traffic and measured delivered background goodput
(Table~\ref{tab:calib}); the maximum, 6.48~Mb/s, is taken as the saturation capacity, giving Low/Medium/High/Saturated $=$
0.65/2.59/5.18/7.77~Mb/s offered. The uplink of C-up is $0.3\times6.48=1.94$~Mb/s.
\begin{table}[t]\centering\caption{Saturation calibration of the 50-node mesh (scenario B, one run per point).}\label{tab:calib}
\resizebox{\columnwidth}{!}{\begin{tabular}{@{}rrrr@{}}\toprule
offered (Mb/s) & carried (Mb/s) & bg goodput (Mb/s) & bg PDR (\%)\\\midrule
1 & 1.04 & 0.97 & 93.0\\
2 & 1.95 & 1.73 & 88.7\\
3 & 2.97 & 2.29 & 77.1\\
4 & 4.30 & 3.15 & 73.4\\
6 & 6.11 & 5.09 & 83.4\\
8 & 8.28 & 6.09 & 73.5\\
10 & 9.76 & 6.48 & 66.3\\
12 & 11.88 & 6.39 & 53.8\\
\bottomrule\end{tabular}}\end{table}

\subsection{Metrics and statistics}
Per severity class: delivery ratio (PDR), mean/median/95th-percentile delay and jitter (mean absolute difference of consecutive delays) of \emph{delivered} packets,
and deadline success (fraction of \emph{sent} packets delivered within 2~s or 5~s; undelivered counts as a failure). Delay is always read together with PDR. Also: background
goodput (to detect starvation), AODV control packets/s, drops by cause (queue-disc overflow, push-out, MAC queue expiry, MAC retry limit,
no route), hop count of delivered packets, and for the burst the PDR and delay of critical packets.
Each configuration is simulated with independent seeds (\texttt{RngRun}) and the \emph{same workload file} in all scenarios, so scenarios are compared
pairwise per seed. We report means with 95\% Student-$t$ confidence intervals over runs, paired two-sided Wilcoxon signed-rank tests (ties/zero differences
handled; if all paired differences are zero the result is reported as ``no difference''), the matched-pairs rank-biserial correlation as effect size,
and Holm correction within each metric family. We make no significance claim without a test.

\subsection{Deviations from the original plan}
Our plan called for 100 nodes, 200~s windows and 30 runs. On the available compute (a loaded 100-node, 70~s mesh run costs about 3 CPU-minutes) the full
grid would take days, so: the default scale is $n=50$ with 8 runs per configuration (6/5/3/3 for ablation/robustness/$n{=}100$/density), 60~s windows, 8 APs instead of 4--6
(Section~\ref{sec:validation}), TCP and DSDV are not evaluated (code paths exist but were not run), the ``priority route discovery'' option is not implemented, and we use ns-3.44.
With 8 runs and about 23 critical packets per run the study can detect large effects only; Section~\ref{sec:threats} discusses this.

\section{Validation: gates and the artefacts they caught}\label{sec:validation}
We defined sanity gates before running experiments. They passed only after fixing four problems, each of which would have produced a wrong conclusion:
\begin{enumerate}
\item \textbf{Route-reply storm.} With ns-3's AODV defaults the 100-node mesh generated runaway control traffic even at zero background load (about 4\,000 control packets/s; RREQ/RREP/RERR/HELLO counts 8372/7965/2071/4737 in the pilot), about 4 minutes of wall time per 30~s of simulated time, and only 59--83\% of warm-up pings arrived. Setting the RFC~3561 ``destination only'' flag (no replies from intermediate caches) reduced the counts to 1777/35/134/246 in the 18~s check run, cut wall time about 10$\times$ (91.5~s to 9.0~s for the same run) and raised warm-up delivery to 97--98\%.
\item \textbf{Gray zone.} Route discovery broadcasts used the basic rate, so AODV selected links that fixed-MCS3 unicast could not sustain:
zero-load SOS delivery was 87.2\% (MCS3 unicast, basic-rate broadcast; $n=100$, run~1) versus 97.2\% when broadcasts are sent at the data MCS (98.9\% with MCS0 everywhere).
\item \textbf{Coverage hole.} With one AP per floor, 7 of 100 stations never associated (best-AP signal $-82.4$ to $-84.2$~dBm, below ns-3's $-82$~dBm
preamble-detection threshold), handicapping C at zero load; two APs per floor associate 100/100.
\item \textbf{Wired-backhaul artefact.} Modelling the backhaul as one shared 5~ms CSMA bus capped the wired side at about $1/5\,\text{ms}=200$ frames/s because ns-3's channel is
busy for the whole propagation delay: at 4.3~Mb/s offered ($n=50$) scenario C delivered only 2.12~Mb/s with 51.8\% SOS delivery and \emph{no} wireless drops. A switched star
of full-rate links delivers 100\%. This would have wrongly shown infrastructure Wi-Fi collapsing under load.
\end{enumerate}
Remaining gates (results on the final configuration, $n=50$, 3 runs, zero background load, 60~s): (G3) SOS delivery 96.7--97.7\% in the mesh
and 100\% in C; A, B, and the EDCA-only/queue-only variants differ by at most 1 percentage point (about 2 packets in 190), as expected when nothing is congested;
(G4) the workload hash is identical for A, B, C in every run; (G5) 49.7 of 50 warm-up pings reach the desk before measurement starts and 50/50 stations are associated in C;
(G6) delivered SOS take 1--5 hops (mean 2.3), i.e.\ the mesh is genuinely multi-hop; (G7) the queue-disc unit test passes (13/13 checks) and, with priority on,
critical packets are \emph{relayed} through AC\_VO queues at intermediate nodes (28 relayed versus 23 originated in run~1), so the marking survives forwarding.
